\section{Discussion}
\label{sec:discussion}

\subsection{Limitations}
\label{subsec:limitations}

\paragraph{Confident wrongness.} \eps cannot detect a model that is
wrong and confident. If training placed all mass on a deprecated API,
the token probability is sharp and \eps is low; the system emits the
wrong code with \textsc{complete} status. Scenario B is the concrete
example: GPT-4o emitted \texttt{openai.ChatCompletion.create(\ldots)}
--- the v0 API, which raises \texttt{AttributeError} on SDK $\ge 1.0$
--- with no hesitation at that token; \eps fired only because of
adjacent uncertainty in surrounding parameter design. Wang et al.~%
\cite{wang2025deprecated} report this failure mode at scale: $25$--$38\%$
deprecated-API emission on standard prompts across seven LLMs, driven
by confident predictions from training data that predates the
deprecation. This is a fundamental limit of any generation-time signal
and a primary motivation for pairing \eps with static analysis
(\S\ref{subsec:static-cov}) rather than replacing it.

\paragraph{Confidently-resolved domains.} Prompts where every model
has converged on a strong training-data prior produce $\eps = 0.000$
regardless of correctness. Current-UTC datetime retrieval and
HTTP-GET-to-JSON fall here across all four models. For stable
widely-deployed APIs this is desirable; for deprecated-but-%
confidently-recalled APIs it is Scenario B's failure mode.

\paragraph{Small-sample calibration.} Main calibration is 30 prompts
per model; Wilson intervals on HIGH detection are $\pm 17$pp at
$n=20$. We claim recall-first design plus an empirical miss rate of
zero on the sets evaluated, not a proof of perfect recall. The \papg
comparison in \S\ref{subsec:papg-comp} uses the same 30 prompts,
which bounds all statistical claims in that subsection to the same
small-sample regime; cross-validation would strengthen all claims on
both sides.

\paragraph{\papg comparison scope.} The \papg evaluation uses the
calibration benchmark (30 prompts per model, 120 total) and has not
been extended to the 201-entry production-plus-scenario set used for
the review-loop results. Whether \papg's better precision on some
models translates to better end-to-end precision in a full two-stage
system, and whether its lack of token-level attribution actually
hurts reviewer performance in practice, are open empirical questions.
The signal-level comparison we report should not be read as a
system-level comparison; a full \papg-routed review-loop experiment
is future work.

\paragraph{Multi-sample sensitivity.} The multi-sample comparison uses
$N=5$ at $T=0.7$. Higher $N$ or higher $T$ might recover additional
behavioural diversity on some prompts --- a $T=1.0$ run would nearly
double the sampling probability of the minority branch on a $0.52/0.47$
token --- though the structural argument does not depend on parameter
choice: \eps is read at $T=0$ where the model's internal distribution
is exposed directly, while any sampling-based method is limited to
what that distribution produces as behaviour. A higher-$T$
configuration would also shift the baseline against which
\emph{generation-time} cost is compared, since production code
generation is typically run at $T \le 0.3$ or deterministically.

\paragraph{Type 3 floor.} Implementation-approach uncertainty on
simple prompts (brute vs.\ hash two-sum) legitimately raises \eps on
correct code. AST filtering does not remove it; the $\sim 50\%$ LOW
FP is predominantly Type 3 and unlikely to shrink without sacrificing
recall.

\paragraph{LLM-judged ground truth.}\label{subsec:gt-disclosure}
LOW ground truth is \texttt{exec()}-verified against assertions
covering return type, edge cases, and correctness. HIGH ground truth
is LLM-judged via GPT-4o. Circularity exists where the evaluated
model shares a family with the judge (GPT-4o-mini judged by GPT-4o);
results for that model should be read with that caveat. We disclose
rather than overclaim human verification we did not perform.

\paragraph{Hyperparameters.} The cascaded tier algorithm (six checks,
three thresholds) was calibrated on GPT-4o; other models inherit
these. The binary \textsc{flagged+} gate is robust; internal tier
placement is a refinement on top.

\subsection{\textsc{paused} tier and delivery semantics}
\label{subsec:paused}

\S\ref{subsec:tiers} defined \textsc{paused} as $\eps \ge 0.95$.
Whether \textsc{paused} actually holds code delivery until the review
completes, or simply runs the review loop like \textsc{flagged}, is
an implementation choice we do not prescribe. In our four-model
scenario sweep no HIGH prompt crossed the \textsc{paused} threshold,
so the question did not become binding. Production use would likely
default to the non-blocking treatment (same routing as
\textsc{flagged}) while exposing \textsc{paused} as a higher-urgency
signal in whatever UI surfaces reviewer output.

\subsection{Context learning as a microcosm}
\label{subsec:context}

Injecting learned false-positive patterns into the reviewer's context
(the configuration evaluated in \S\ref{subsec:review-loop}) raised
clearance from $53.2\%$ to $68.7\%$, and late false positives fell from
$4$ to $0$. This is in-context improvement on a fixed model; training-%
time improvement would be expected to compound with it. The same
mechanism --- priming the reviewer with model-specific priors from a
prior clean run --- is a candidate for the same kind of systematic
study that has emerged for in-context learning more broadly.

\subsection{Future work}
\label{subsec:future}

\textbf{Pipeline propagation:} upstream \eps should propagate as
prior uncertainty into downstream generations, modulating sampling
temperature or review threshold. \textbf{Open-source models:} vLLM
exposes full top-$k$; a Llama/Qwen/Mistral calibration run would test
whether the recall-first boundary holds outside OpenAI. \textbf{Per-%
token intervention:} surface the high-\eps token as an inline editor
marker, or trigger targeted re-decoding with lower temperature on
just the uncertain span (AdaDec-style, but at review time).
\textbf{Cross-decision consistency:} Scenario D's mismatches argue
for AST-level consistency checks over \eps (e.g., \texttt{async def}
$+$ blocking HTTP client). \textbf{Import-gated activation:}
restricting \eps scoring to functions that import external libraries
(i.e., skipping pure-logic functions) would eliminate the Type 3
false-positive floor at negligible cost, since production API-%
integration code almost always begins with an import.
\textbf{Combined \papg + \eps routing:} \papg's function-level
specificity could be combined with \eps's token localization --- a
two-signal gate that fires when both agree may recover precision
without sacrificing recall. \textbf{Full-scale \papg comparison:}
extending the \papg evaluation to the 201-entry review-loop set
would convert the signal-level comparison of \S\ref{subsec:papg-comp}
into a system-level one.
